\documentclass[12pt]{article}

\usepackage[utf8]{inputenc}
\usepackage{amsmath}
\usepackage{amssymb}
\usepackage{graphicx}
\usepackage{bbm}
\usepackage{jheppub}

\newcommand{\Ident}{\mathbbm{1}}
\newcommand{\svx}{\mathsf{x}}
\newcommand{\hws}{\mathrm{hws}}
\newcommand{\gl}{\mathfrak{gl}}
\newcommand{\so}{\mathfrak{so}}
\newcommand{\su}{\mathfrak{su}}
\newcommand{\Sym}{\mathrm{Sym}}
\newcommand{\Gr}{\mathrm{Gr}}
\newcommand{\wt}{\mathrm{wt}}

\title{Gaudin CG}

\author{Paul Ryan}

\emailAdd{paul.ryan@desy.de}

\abstract{This is the start of the gaudin-CG document.}

\begin{document}

\maketitle

\section{Introduction}

This is the start of the gaudin-CG document.

\section{Rectangular $\gl(4)$ representations and the abelian nilradical}
\label{sec:rect-nilradical}

This section sets up the representation-theoretic structure underlying the
Clebsch--Gordan (CG) problem for rectangular $\gl(4)$ representations
$\lambda=[S,S,0,0]$. The upshot is that for these representations, and only for
these, the highest-weight covector is annihilated by the entire off-diagonal part
of a Levi subalgebra $\gl(2)\oplus\gl(2)$, so that the whole dual module is
generated by four \emph{commuting} raising operators. This makes the CG problem
combinatorially explicit: monomials in those four operators label states, the
weight of the ket truncates the relevant exponential to a finite sum, and the
complete set of relations between the labels can be classified in closed
form.\footnote{All statements below marked as verified were checked by explicit
matrix computation in the Gelfand--Tsetlin basis in the notebook
\texttt{Paul/Mathematica/Gaudin/Clean/CGGenOp.wb}; the verification ranges quoted
are the ones actually run.}

\subsection{Setup and the extra annihilation conditions}

Let $E_{ij}$, $i,j=1,\dots,4$, denote the $\gl(4)$ generators, with
\[
[E_{ij},E_{kl}] = \delta_{jk}E_{il}-\delta_{li}E_{kj},
\]
acting on the finite-dimensional irreducible representation $V_\lambda$ of
highest weight $\lambda=(\lambda_1,\lambda_2,\lambda_3,\lambda_4)$, realized on
the Gelfand--Tsetlin basis. Let $|\hws\rangle$ be the highest-weight vector,
annihilated by $E_{jk}$ for $j<k$, and $\langle\hws|$ the dual covector. For any
$\lambda$ one has $\langle\hws|E_{jk}=0$ for $j>k$, since $\langle\hws|E_{jk}$
would pair with kets of weight $\lambda+e_j-e_k \not\leq \lambda$.

Now specialize to the rectangular case $\lambda=[S,S,0,0]$. The weights
\[
\lambda-e_1+e_2=(S-1,S+1,0,0),\qquad
\lambda-e_3+e_4=(S,S,-1,1)
\]
do not occur in $V_\lambda$, because their dominant rearrangements
$(S+1,S-1,0,0)$ and $(S,S,1,-1)$ are not $\leq\lambda$. Hence
\begin{equation}
\langle\hws|E_{12}=\langle\hws|E_{34}=0
\label{eq:extra-annihilation}
\end{equation}
on top of the generic conditions $\langle\hws|E_{21}=\langle\hws|E_{43}=0$. Thus
$\langle\hws|$ is annihilated by the complete off-diagonal part of the Levi
subalgebra $\gl(2)\oplus\gl(2)\subset\gl(4)$ acting on the index blocks
$\{1,2\}$ and $\{3,4\}$: it is an extreme vector for the corresponding parabolic.
Conditions \eqref{eq:extra-annihilation} fail for non-rectangular
representations; they were checked to hold for $S=1,2,3$ and to fail for
$\lambda=(3,2,1,0)$ and $\lambda=(2,1,1,0)$.

The surviving generators are those of the nilradical
\begin{equation}
\mathfrak{n}^+=\mathrm{span}\{E_{13},E_{14},E_{23},E_{24}\}\;\cong\;\mathbb{C}^2\otimes\mathbb{C}^2 ,
\end{equation}
and this nilradical is \emph{abelian}: with $i,k\in\{1,2\}$ and $j,l\in\{3,4\}$
drawn from disjoint index sets, every Kronecker delta in the commutator vanishes,
so
\begin{equation}
[E_{ij},E_{kl}]=0,\qquad i,k\in\{1,2\},\ j,l\in\{3,4\}.
\label{eq:abelian}
\end{equation}
Consequently the four operators may be multiplied in any order, and any product
of the corresponding exponentials collapses to a single exponential,
\begin{equation}
e^{a_{13}E_{13}}e^{a_{14}E_{14}}e^{a_{23}E_{23}}e^{a_{24}E_{24}}
= \exp\Big[\sum_{i\in\{1,2\}}\sum_{j\in\{3,4\}} a_{ij}E_{ij}\Big],
\label{eq:merge}
\end{equation}
verified as an exact matrix identity for $\lambda=(3,2,1,0)$. Writing
$x_{ij}:=E_{ij}$, the map
\begin{equation}
\mathbb{C}[x_{13},x_{14},x_{23},x_{24}]\;\longrightarrow\;V_\lambda^*,\qquad
P\;\longmapsto\;\langle\hws|\,P ,
\label{eq:poly-map}
\end{equation}
is surjective: the rank of the set of monomial covectors equals $\dim V_\lambda$
exactly, checked for $S=1,2,3$ ($\dim=6,20,50$). Four commuting operators
therefore generate the entire dual module.

\subsection{Weight grading of the monomial expansion}

By \eqref{eq:abelian} and \eqref{eq:merge} the expansion of the exponential is a
plain multinomial with no ordering ambiguity,
\begin{equation}
\langle\hws|\,e^{\sum a_{ij}E_{ij}}
=\sum_{n_{13},n_{14},n_{23},n_{24}\geq0}
\ \prod_{ij}\frac{a_{ij}^{\,n_{ij}}}{n_{ij}!}\
\langle\hws|E_{13}^{n_{13}}E_{14}^{n_{14}}E_{23}^{n_{23}}E_{24}^{n_{24}} .
\label{eq:expansion}
\end{equation}
The covector $\langle\hws|E_{13}$ pairs nontrivially only with kets of weight
$\lambda-e_1+e_3$, since $E_{13}$ must raise such a ket back to $|\hws\rangle$.
Labelling covectors by the ket weight they pair with, the monomial
$n=(n_{13},n_{14},n_{23},n_{24})$ therefore carries weight
\begin{equation}
\nu(n)=\lambda-(n_{13}+n_{14})e_1-(n_{23}+n_{24})e_2+(n_{13}+n_{23})e_3+(n_{14}+n_{24})e_4 .
\label{eq:weight}
\end{equation}
Only the row and column sums of the $2\times2$ exponent table
$\binom{n_{13}\ n_{14}}{n_{23}\ n_{24}}$ enter, that is the margins
\begin{equation}
m_1=n_{13}+n_{14},\quad m_2=n_{23}+n_{24},\quad
m_3=n_{13}+n_{23},\quad m_4=n_{14}+n_{24},
\label{eq:margins}
\end{equation}
subject to $m_1+m_2=m_3+m_4$, which is automatic since each $E_{ij}$ preserves
the total weight. A weight sector is thus degenerate: fixing $m$ leaves the
one-parameter family
\begin{equation}
\begin{gathered}
n_{13}=t,\quad n_{14}=m_1-t,\quad n_{23}=m_3-t,\quad n_{24}=m_2-m_3+t,\\
\max(0,m_3-m_2)\leq t\leq\min(m_1,m_3),
\end{gathered}
\label{eq:tfamily}
\end{equation}
consisting of exactly $\min(m_1,m_2,m_3,m_4)+1$ monomials. The grading
\eqref{eq:weight} was verified for $\lambda=(3,2,1,0)$ by checking that every
monomial covector of total degree $\leq6$ is supported only on Gelfand--Tsetlin
patterns of weight $\nu(n)$.

\subsection{Overlaps with a ket of fixed weight}

Let $|\Psi\rangle\in V_\lambda$ have definite weight $\mu$. Only the monomials
with $\nu(n)=\mu$ contribute to $\langle\hws|e^{\sum a_{ij}E_{ij}}|\Psi\rangle$;
all others land in a different weight space and pair to zero. Inverting
\eqref{eq:weight} fixes the margins completely,
\begin{equation}
m_1=\lambda_1-\mu_1,\qquad m_2=\lambda_2-\mu_2,\qquad
m_3=\mu_3-\lambda_3,\qquad m_4=\mu_4-\lambda_4 ,
\label{eq:forced-margins}
\end{equation}
so that the overlap vanishes identically unless all four $m_i\geq0$ (and
$\sum_i\mu_i=\sum_i\lambda_i$). When these conditions hold, the infinite series
\eqref{eq:expansion} collapses to the single family \eqref{eq:tfamily},
\begin{equation}
\boxed{\ \
\begin{aligned}
&\langle\hws|e^{\sum a_{ij}E_{ij}}|\Psi\rangle\\[2pt]
&\qquad=\sum_{t=\max(0,m_3-m_2)}^{\min(m_1,m_3)}
\frac{a_{13}^{\,t}\,a_{14}^{\,m_1-t}\,a_{23}^{\,m_3-t}\,a_{24}^{\,m_2-m_3+t}}
{t!\,(m_1-t)!\,(m_3-t)!\,(m_2-m_3+t)!}\\[2pt]
&\qquad\qquad\qquad\qquad\times\
\langle\hws|E_{13}^{t}E_{14}^{m_1-t}E_{23}^{m_3-t}E_{24}^{m_2-m_3+t}|\Psi\rangle
\end{aligned}
\ \ }
\label{eq:overlap}
\end{equation}
a finite sum of exactly $\min(m)+1$ terms, homogeneous of degree $m_1+m_2$ in the
$a_{ij}$. No truncation argument is needed: fixing the target weight does all the
work. In particular the overlap reduces to a single matrix element whenever
$\min(m)=0$.

Formula \eqref{eq:overlap} was verified against the brute-force matrix
exponential for $\lambda=(3,2,1,0)$ at all $38$ weights occurring in the
representation, with $|\Psi\rangle$ a random integer combination of the
Gelfand--Tsetlin vectors at each weight; the difference vanishes exactly,
including at the $12$ weights where the right-hand side is predicted to be
empty. It was checked again for the rectangular case $\lambda=(4,4,0,0)$,
$\mu=(2,2,2,2)$.

\subsection{The middle weight of a rectangular representation}

The case relevant to the CG application is $\lambda=[S,S,0,0]$ with
$|\Psi\rangle$ at the middle weight $\mu=(\tfrac{S}{2},\tfrac{S}{2},\tfrac{S}{2},\tfrac{S}{2})$
for $S$ even, where \eqref{eq:forced-margins} gives
$m=(\tfrac{S}{2},\tfrac{S}{2},\tfrac{S}{2},\tfrac{S}{2})$. Then $m_3-m_2=0$ and
$m_1=m_3$, so $t$ runs over $0,\dots,\tfrac{S}{2}$ and the exponents are locked
into the anti-diagonally symmetric pattern
$n=(t,\tfrac{S}{2}-t,\tfrac{S}{2}-t,t)$. The overlap becomes
\begin{equation}
\begin{gathered}
\langle\hws|e^{\sum a_{ij}E_{ij}}|\Psi\rangle
=\sum_{t=0}^{S/2}\frac{x^{t}\,y^{\,S/2-t}}{\big(t!\,(\tfrac{S}{2}-t)!\big)^{2}}\
\langle\hws|E_{13}^{t}E_{14}^{S/2-t}E_{23}^{S/2-t}E_{24}^{t}|\Psi\rangle,\\
x:=a_{13}a_{24},\qquad y:=a_{14}a_{23} .
\end{gathered}
\label{eq:middle-weight}
\end{equation}
Because $n_{13}=n_{24}$ and $n_{14}=n_{23}$ are forced, the overlap depends on
the four twist parameters only through the two combinations $x$ and $y$: it is a
homogeneous polynomial of degree $S/2$ in $(x,y)$, whose $S/2+1$ coefficients are
precisely the individual matrix elements. Reading off powers of $x/y$ therefore
separates all of them.

For $S=16$, $\mu=(8,8,8,8)$, the margins are $m=(8,8,8,8)$ and exactly nine
monomials contribute, all of total degree $16$:
\[
\begin{array}{c|c}
(n_{13},n_{14},n_{23},n_{24}) & \text{coefficient} \\ \hline
(0,8,8,0) & a_{14}^{8}a_{23}^{8}/(8!)^2 \\
(1,7,7,1) & a_{13}a_{14}^{7}a_{23}^{7}a_{24}/(7!)^2 \\
(2,6,6,2) & a_{13}^{2}a_{14}^{6}a_{23}^{6}a_{24}^{2}/(2!\,6!)^2 \\
\ \vdots\ & \vdots \\
(4,4,4,4) & a_{13}^{4}a_{14}^{4}a_{23}^{4}a_{24}^{4}/(4!\,4!)^2 \\
\ \vdots\ & \vdots \\
(8,0,0,8) & a_{13}^{8}a_{24}^{8}/(8!)^2
\end{array}
\]
in agreement with \eqref{eq:middle-weight}. As shown in
section~\ref{sec:relations}, this is exactly the borderline case in which the
nine monomials are linearly independent, so no relations contaminate the
extraction of the matrix elements.

\section{Classification of the relations}
\label{sec:relations}

The map \eqref{eq:poly-map} is surjective but not injective: the monomial
covectors are an overcomplete spanning set. This section identifies the kernel
$I=\ker$ completely.

\subsection{Weight multiplicities}

Within the weight sector of margins $m$ there are $\min(m)+1$ monomials, but the
multiplicity of the weight $\mu$ in $V_{[S,S,0,0]}$ is only
\begin{equation}
\dim V_{[S,S,0,0]}[\mu] \;=\; \min\big(\min(m)+1,\ S+1-\max(m)\big),
\label{eq:multiplicity}
\end{equation}
with $m$ as in \eqref{eq:forced-margins}, verified against the
Gelfand--Tsetlin count at every weight for $S=1,2,3,4$ ($85$ weights at $S=4$).
The two counts agree if and only if $\min(m)+\max(m)\leq S$; otherwise the
monomials of the sector are linearly dependent, the number of independent
relations in that sector being
\begin{equation}
\max\big(0,\ \min(m)+\max(m)-S\big).
\label{eq:sector-deficit}
\end{equation}
The middle weight $m=(\tfrac{S}{2},\dots,\tfrac{S}{2})$ saturates the bound,
$\min(m)+\max(m)=S$, which is why \eqref{eq:middle-weight} is relation-free.

\subsection{The relations are the rank-one degenerations}

The ideal $I$ is generated in the single degree $S+1$, by the Cartan component
\begin{equation}
\Sym^{S+1}\mathbb{C}^2\otimes\Sym^{S+1}\mathbb{C}^2\ \subset\ \Sym^{S+1}(\mathbb{C}^2\otimes\mathbb{C}^2),
\end{equation}
of dimension $(S+2)^2$. Equivalently, and far more usefully, the relations are
exactly the statement that a \emph{rank-one} twist matrix $a_{ij}=u_iv_j$ makes
the $(S+1)$-st power annihilate the covector:
\begin{equation}
\boxed{\ \
\langle\hws|\Big(\sum_{i\in\{1,2\}}\sum_{j\in\{3,4\}} u_i v_j\,E_{ij}\Big)^{S+1}=0
\ \ }
\label{eq:rank-one}
\end{equation}
identically in $u_1,u_2,v_3,v_4$, while the $S$-th power does not vanish. Both
statements were verified symbolically in $u,v$ for $S=1,2,3$.

Extracting from \eqref{eq:rank-one} the coefficient of
$u_1^{m_1}u_2^{m_2}v_3^{m_3}v_4^{m_4}$, necessarily with
$m_1+m_2=m_3+m_4=S+1$, produces one relation per weight sector, supported on the
entire $t$-family \eqref{eq:tfamily} of that sector with multinomial weights:
\begin{equation}
R_m:\qquad
\sum_{t}\frac{(S+1)!}{t!\,(m_1-t)!\,(m_3-t)!\,(m_2-m_3+t)!}\
\langle\hws|E_{13}^{t}E_{14}^{m_1-t}E_{23}^{m_3-t}E_{24}^{m_2-m_3+t}=0 .
\label{eq:Rm}
\end{equation}
There are $(S+2)^2$ such relations, labelled by $0\leq m_1,m_3\leq S+1$, which
matches $\dim(\Sym^{S+1}\otimes\Sym^{S+1})$. The coefficients in \eqref{eq:Rm}
are precisely those of the overlap formula \eqref{eq:overlap}: the relation is
the sector sum of \eqref{eq:overlap} at the rank-one locus. All $(S+2)^2$
relations \eqref{eq:Rm} were verified to vanish exactly for $S=1,2,3$, and the
analogous degree-$S$ combinations were checked \emph{not} to vanish, so the
degree $S+1$ is sharp.

No further generators occur. Writing $K_k\subset\Sym^k(\mathbb{C}^4)$ for the
degree-$k$ part of $I$, it was verified for $S=1,2$ that $K_k=0$ for $k\leq S$,
that $\dim K_{S+1}=(S+2)^2$, and that
$\mathrm{rank}\big(x\cdot K_k\big)=\dim K_{k+1}$ for $k\geq S+1$, i.e.\ the ideal
is generated in degree $S+1$ alone.

\subsection{Graded structure}

Let
\begin{equation}
D:=E_{13}E_{24}-E_{14}E_{23}
\end{equation}
denote the determinant of the $2\times2$ matrix of nilradical generators, which
is well defined by \eqref{eq:abelian}. Grading $V_\lambda^*$ by the degree in the
$E_{ij}$, the decomposition into irreducibles of the Levi
$\gl(2)\oplus\gl(2)$ is
\begin{equation}
V_{[S,S,0,0]}^*\;\cong\;\bigoplus_{a=0}^{S}\ \bigoplus_{j=0}^{S-a}\
D^{\,j}\cdot\big(\Sym^a\mathbb{C}^2\otimes\Sym^a\mathbb{C}^2\big),
\label{eq:graded}
\end{equation}
so that the degree-$k$ subspace has dimension
$\sum_{j}(k-2j+1)^2$ restricted to $0\leq k-2j\leq S$ and $j\leq S-(k-2j)$, and
\begin{equation}
\dim V_{[S,S,0,0]}=\sum_{a=0}^{S}(a+1)^2(S+1-a)=\frac{(S+1)(S+2)^2(S+3)}{12}.
\label{eq:dim}
\end{equation}
The graded dimensions predicted by \eqref{eq:graded} were confirmed exactly for
$S=1,2,3$, namely $\{1,4,1\}$, $\{1,4,10,4,1\}$ and $\{1,4,10,20,10,4,1\}$; note
that they are palindromic under $k\leftrightarrow 2S-k$. It follows from
\eqref{eq:graded} that $\langle\hws|D^{S+1}=0$, verified directly for $S=1,2,3$,
whereas $\langle\hws|D\neq0$. We emphasize that $D^{S+1}$ is a
\emph{consequence} of the relations \eqref{eq:Rm} and not a generator of $I$: it
sits in degree $2S+2$, well above the generation degree $S+1$.

\subsection{A basis: Pl\"ucker and standard-monomial labels}

The obstruction to using monomials in the four $E_{ij}$ as a basis is that $D$ is
not expressible as such a monomial modulo $I$. Adjoining it as a fifth generator
resolves the problem. On the big cell of the Grassmannian $\Gr(2,4)$, the
Pl\"ucker coordinates of the matrix
$\left(\begin{smallmatrix}1&0&x_{13}&x_{14}\\0&1&x_{23}&x_{24}\end{smallmatrix}\right)$
are
\begin{equation}
p_{12}=1,\qquad p_{13}=E_{13},\quad p_{14}=E_{14},\quad
p_{23}=E_{23},\quad p_{24}=E_{24},\qquad p_{34}=D ,
\end{equation}
and $V_{[S,S,0,0]}^*$ is the space of degree-$S$ monomials in the six $p$'s
modulo the single Klein quadric
\begin{equation}
p_{12}p_{34}-p_{13}p_{24}+p_{14}p_{23}=0 .
\label{eq:klein}
\end{equation}
Standard monomial theory then instructs us to eliminate the unique incomparable
pair $p_{14}p_{23}$, giving the basis
\begin{equation}
\boxed{\ \
\Big\{\ \langle\hws|\,E_{13}^{e_{13}}E_{14}^{e_{14}}E_{23}^{e_{23}}E_{24}^{e_{24}}D^{e_{34}}
\ :\ \textstyle\sum e\leq S,\quad \min(e_{14},e_{23})=0\ \Big\}
\ \ }
\label{eq:basis}
\end{equation}
the slack $e_{12}=S-\sum e$ filling the remaining columns. These labels are in
bijection with the semistandard Young tableaux of shape $(S,S)$ on the alphabet
$\{1,2,3,4\}$, as they must be. That \eqref{eq:basis} is an exact basis --- the
number of labels equals the rank of the corresponding covectors equals
$\dim V_{[S,S,0,0]}$ --- was verified for $S=1,2,3$. Five integers subject to two
transparent constraints thus replace four integers subject to the hidden
relations \eqref{eq:Rm}, at no cost in convention-dependence: no
Gelfand--Tsetlin patterns and no phase choices enter.

\subsection{How rich are the rectangular representations?}

It is worth quantifying how much structure \eqref{eq:graded} actually contains.
Compare the symmetric powers $\lambda=[S,0,0,0]$: there the Levi is
$\gl(1)\oplus\gl(3)$, the nilradical is spanned by the \emph{three} commuting
operators $E_{12},E_{13},E_{14}$, and the monomials of degree $\leq S$ already
form a basis. The ideal is then simply everything of degree $>S$, there are no
relations at all, and every weight multiplicity is $1$. This was verified for
$S=2,3,4$, where the monomial count, the rank and $\binom{S+3}{3}$ all coincide,
and where $E_{23},E_{24},E_{34}$ additionally annihilate the covector. Symmetric
powers are in this sense a truncated free polynomial ring.

Rectangular representations are strictly richer: weight multiplicities reach
$\lfloor S/2\rfloor+1$, attained at the middle weight, since
\eqref{eq:multiplicity} is maximized over $m$ at $m_1=m_2=m_3=m_4=\lfloor
S/2\rfloor$ (verified for $S=1,\dots,5$); there is a genuine relation, and there is a
nontrivial filtration by powers of $D$. But the additional richness amounts to
exactly one extra integer, the exponent $e_{34}$ of $D$, because $\Gr(2,4)$ is a
single quadric \eqref{eq:klein} in $\mathbb{P}^5$. The invariant way to say this
uses $\su(4)\cong\so(6)$:
\begin{equation}
[S,S,0,0]\ \cong\ \mathcal{H}^S(\mathbb{C}^6),
\end{equation}
the traceless symmetric degree-$S$ tensors on $\mathbb{C}^6$, i.e.\ the spherical
harmonics of $\so(6)$, with $D$ playing the role of the invariant quadratic form.
The corresponding dimension identity
\begin{equation}
\frac{(S+1)(S+2)^2(S+3)}{12}=\binom{S+5}{5}-\binom{S+3}{5}
\end{equation}
was checked for $S=1,\dots,6$. Rectangular representations are therefore as rich
as spherical harmonics --- one step up from free, but governed throughout by a
single quadric, which is what makes the closed-form CG generating function
possible. The genuinely hard case remains generic
$\lambda=[\lambda_1,\lambda_2,\lambda_3,\lambda_4]$, where $\langle\hws|$ is no
longer a parabolic extreme vector, the relevant nilradical is not abelian, and
none of the above simplifications apply.

\bibliographystyle{JHEP}
\bibliography{references}

\end{document}
